\section{Partie théorique}

\subsection{Principe de l'actionneur à bobine mobile}

L'actionneur étudié est un transducteur électrodynamique à déflexion angulaire limitée, utilisé en particulier pour le positionnement des têtes de lecture des disques durs. Une bobine plate, solidaire d'un axe rotatif, plonge dans l'entrefer d'un aimant permanent semi-circulaire qui crée un champ radial \(B\) constant. Lorsqu'un courant \(i\) traverse la bobine, la force de Laplace produit un couple moteur \(T_{em} = K_t\,i\). Un ressort de torsion impose une position d'équilibre stable \(\alpha = 0\) en l'absence de courant et limite la course angulaire à \(\pm\Delta\alpha_{\max}\).

\begin{figure}[H]
\centering
\begin{tikzpicture}[>=Stealth, scale=0.92, transform shape,
  every node/.style={font=\small}]
% Repere et rayons
\coordinate (O) at (0,0);
\draw[thick] (O) -- (118:2.4);
\draw[thick] (O) -- (62:2.4);
\draw[densely dashed] (O) -- (90:2.8);
\filldraw[black] (O) circle (1.2pt) node[below=2pt] {\scriptsize axe};
% Aimants et entrefer
\draw[very thick, cyan!70!black] (118:1.55) arc (118:62:1.55);
\draw[very thick, cyan!70!black] (118:2.25) arc (118:62:2.25);
\draw[very thick, cyan!70!black] (118:1.55) -- (118:2.25);
\draw[very thick, cyan!70!black] (62:1.55) -- (62:2.25);
\node[font=\scriptsize, cyan!60!black] at (111:1.9) {N};
\node[font=\scriptsize, cyan!60!black] at (69:1.9) {S};
% Bobine mobile
\draw[very thick, orange!85!black] (104:1.72) arc (104:76:1.72);
\draw[very thick, orange!85!black] (104:2.08) arc (104:76:2.08);
\draw[very thick, orange!85!black] (104:1.72) -- (104:2.08);
\draw[very thick, orange!85!black] (76:1.72) -- (76:2.08);
\foreach \ang in {80,85,90,95,100} {
  \draw[orange!85!black] (\ang:1.73) -- (\ang:2.07);
}
% Couple
\draw[->, thick, red!85!black] (0,2.75) arc (90:54:2.75)
  node[pos=0.72,above right] {\(F_{l}\)};
\node[font=\scriptsize] at (1.15,1.25) {\(R_i\)};
\node[font=\scriptsize] at (1.45,1.95) {\(R_e\)};
\draw[<->, thick] (-2.8,1.55) -- (-2.8,2.25) node[midway,left] {\(\ell=R_e-R_i\)};
% Equations
\node[align=left, anchor=west] at (3.1,2.3) {\(F_L = B\,\ell\,i = B(R_e-R_i)i\)};
\node[align=left, anchor=west] at (3.1,1.25)
{\(T_{em}=2\frac{R_e+R_i}{2}F_LN\)};
\node[align=left, anchor=west] at (3.1,0.35)
{\(T_{em}=B(R_e^2-R_i^2)N\,i=K_t i\)};
\node[align=left, anchor=west] at (3.1,-0.65)
{\(u_{im}=K_u\Omega,\qquad K_u=K_t\)};
\end{tikzpicture}
\caption{Principe électromagnétique de l'actionneur à bobine mobile.}
\label{fig:actionneur_principe}
\end{figure}

\begin{figure}[H]
\centering
\begin{tikzpicture}[>=Stealth, node distance=1.5cm]
% Bloc électrique
\draw[thick, rounded corners=4pt] (0,0) rectangle (3.2,2.2);
\node[align=center] at (1.6,1.5) {\small Circuit};
\node[align=center] at (1.6,0.9) {\small électrique};
\node[font=\scriptsize] at (1.6,0.35) {\(R,\, L\)};
% Bloc mécanique
\draw[thick, rounded corners=4pt] (5.8,0) rectangle (9,2.2);
\node[align=center] at (7.4,1.5) {\small Système};
\node[align=center] at (7.4,0.9) {\small mécanique};
\node[font=\scriptsize] at (7.4,0.35) {\(J,\, C_v,\, K_r\)};
% Flèches et couplage
\draw[->] (-1.2,1.1) -- (0,1.1) node[midway,above]{\small \(u\)};
\draw[->] (3.2,1.4) -- (5.8,1.4) node[midway,above]{\small \(K_t\,i\)};
\draw[->] (5.8,0.8) -- (3.2,0.8) node[midway,below]{\small \(K_u\,\dot{\alpha}\)};
\draw[->] (9,1.1) -- (10.2,1.1) node[midway,above]{\small \(\alpha\)};
\end{tikzpicture}
\caption{Schéma de principe : couplage électro-mécanique de l'actionneur}
\label{fig:schema_bloc}
\end{figure}

\subsection{Couple électromagnétique et tension induite}

\subsubsection{Couple électromagnétique}

La force de Laplace exercée sur un conducteur de longueur \(\ell\) parcouru par un courant \(i\) dans un champ magnétique \(B\) s'écrit \(F = B\,\ell\,i\). Pour la bobine étudiée, la longueur active radiale d'un côté de spire vaut \(\ell = R_e - R_i\), où \(R_e\) et \(R_i\) sont les rayons extérieur et intérieur de la bobine. Pour \(N\) spires :
\begin{equation}
F = B\,(R_e - R_i)\,N\,i.
\end{equation}
Chaque spire possède deux côtés actifs diamétralement opposés, situés au rayon moyen \(r = (R_e + R_i)/2\). Le couple électromagnétique total est donc :
\begin{equation}
T_{em} = 2\,r\,F = 2\cdot\frac{R_e + R_i}{2}\cdot B\,(R_e - R_i)\,N\,i = B\,(R_e^2 - R_i^2)\,N\,i.
\end{equation}
On en déduit la constante de couple :
\begin{equation}
K_t = B\,(R_e^2 - R_i^2)\,N.
\label{eq:Kt_geom}
\end{equation}

\subsubsection{Tension induite}

Lorsque la bobine se déplace à la vitesse angulaire \(\dot\alpha\) dans le champ \(B\), la vitesse linéaire d'un côté actif est \(v = r\,\dot\alpha\). La tension induite par côté de spire vaut \(u_i = B\,\ell\,v\). En sommant sur les deux côtés actifs et les \(N\) spires :
\begin{equation}
u_{im} = 2\,B\,(R_e - R_i)\cdot\frac{R_e + R_i}{2}\,N\,\dot\alpha = B\,(R_e^2 - R_i^2)\,N\,\dot\alpha.
\end{equation}
On identifie ainsi la constante de tension induite :
\begin{equation}
K_u = B\,(R_e^2 - R_i^2)\,N.
\label{eq:Ku_geom}
\end{equation}

\paragraph{Identité \(K_u = K_t\) en unités SI.}
Les expressions \eqref{eq:Kt_geom} et \eqref{eq:Ku_geom} montrent que les deux constantes sont géométriquement identiques :
\begin{equation}
\boxed{\,K_u = K_t \equiv K = B\,(R_e^2 - R_i^2)\,N.\,}
\label{eq:Ku_eq_Kt}
\end{equation}
Ce résultat est cohérent avec le bilan énergétique \(P_{m\acute{e}c} = T_{em}\,\dot\alpha = K_t\,i\,\dot\alpha\) et \(P_{\acute{e}lec} = u_{im}\,i = K_u\,\dot\alpha\,i\), qui impose \(K_u = K_t\) en unités SI cohérentes. Cette identité sera utilisée pour la suite, et \eqref{eq:Ku_geom} permettra d'extraire le nombre de spires \(N\) à partir de \(K_u\) mesuré.

\subsection{Équations de mouvement}

\subsubsection{Équation électrique}

La relation électrique peut se voir comme une loi d'Ohm généralisée :
\[
u(t) = R\,i(t) + \frac{d\psi}{dt}, \qquad \psi = L\,i(t) + \psi_m(\alpha).
\]
Le terme \(d\psi_m/dt\) correspond à la tension induite de mouvement. Comme \(\dot\alpha = \Omega\), il vaut \(K_u\,\dot{\alpha}(t)\).
La loi des mailles sur le circuit de la bobine donne :
\begin{equation}
u(t) \;=\; R\,i(t) \;+\; L\,\frac{di}{dt} \;+\; K_u\,\dot{\alpha}(t),
\label{eq:elec}
\end{equation}
où \(K_u\,\dot{\alpha}\) est la force contre-électromotrice (f.c.e.m.) induite par la rotation de la bobine dans le champ.

\subsubsection{Équation mécanique}

Le principe fondamental de la dynamique en rotation s'écrit :
\begin{equation}
J\,\ddot{\alpha}(t) \;=\; K_t\,i(t) \;-\; C_v\,\dot{\alpha}(t) \;-\; K_r\,\alpha(t).
\label{eq:meca}
\end{equation}

\subsubsection{Constantes de temps caractéristiques}

Le système comporte deux dynamiques de natures différentes :
\begin{itemize}
    \item une dynamique \emph{électrique} de constante de temps \(\tau_{el} = L/R\) ;
    \item une dynamique \emph{mécanique} caractérisée par \(\tau_{m\acute{e}c} \sim 1/\delta\) avec \(\delta = C_v/(2J)\), beaucoup plus lente.
\end{itemize}
Sous l'hypothèse usuelle \(\tau_{el} \ll \tau_{m\acute{e}c}\), le terme \(L\,di/dt\) peut être négligé dans \eqref{eq:elec} pour l'analyse mécanique, ce qui simplifie EX2 et EX3 ci-dessous. Ce point sera vérifié numériquement à la section~\ref{sec:synthese}.

\subsection{Stratégie d'identification}
\label{sec:strat_ident}

Les six paramètres sont identifiés séquentiellement à partir de trois essais.

\subsubsection{EX1 -- Réponse indicielle bloquée : \texorpdfstring{\(R\)}{R} et \texorpdfstring{\(L\)}{L}}

L'actionneur est maintenu mécaniquement (\(\dot{\alpha} = 0\)). La f.c.e.m.\ est alors nulle et \eqref{eq:elec} se réduit à un simple circuit RL :
\[
u(t) \;=\; R\,i(t) \;+\; L\,\frac{di}{dt}.
\]
Dans le domaine de Laplace, pour conditions initiales nulles :
\[
G_{el}(s) = \frac{I(s)}{U(s)} = \frac{1}{R + Ls}
= \frac{1}{R}\,\frac{1}{1 + sL/R}.
\]
Le gain statique vaut donc \(1/R\), et la constante de temps vaut \(\tau_{el}=L/R\).
La réponse à un échelon \(u_0\) est :
\begin{equation}
i(t) \;=\; \frac{u_0}{R}\,\bigl(1 - e^{-t/\tau_{el}}\bigr), \qquad \tau_{el} = \frac{L}{R}.
\label{eq:i_RL}
\end{equation}
L'état stationnaire fournit \(R = \Delta u/\Delta i_{\infty}\) et la constante de temps fournit \(L = R\,\tau_{el}\).

\subsubsection{EX2 -- Bobine ouverte : \texorpdfstring{\(\delta\)}{d}, \texorpdfstring{\(\omega_0\)}{w0} et \texorpdfstring{\(K_u\)}{Ku}}

En circuit ouvert, \(i = 0\) et le couple électromagnétique s'annule. L'équation mécanique \eqref{eq:meca} devient :
\begin{equation}
J\,\dot\Omega(t) \;=\; -\,K_r\,\alpha(t) \;-\; C_v\,\Omega(t),
\label{eq:meca_libre_omega1}
\end{equation}
avec \(\Omega = \dot\alpha\). En dérivant \eqref{eq:meca_libre_omega1} par rapport au temps et en utilisant \(\dot\alpha = \Omega\), on élimine la position et on obtient une équation différentielle du second ordre directement en \(\Omega(t)\) :
\begin{equation}
J\,\ddot\Omega + C_v\,\dot\Omega + K_r\,\Omega = 0
\;\Longleftrightarrow\;
\ddot\Omega + 2\delta\,\dot\Omega + \omega_n^2\,\Omega = 0,
\label{eq:meca_libre}
\end{equation}
avec, par identification :
\begin{equation}
\delta = \frac{C_v}{2J}, \qquad \omega_n^2 = \frac{K_r}{J}.
\end{equation}

\paragraph{Solution avec conditions initiales.}
Pour \(\omega_n > \delta\) (régime pseudo-périodique), la solution générale s'écrit (annexe A.1) :
\begin{equation}
\Omega(t) = e^{-\delta t}\bigl[C_1\cos(\omega_0\,t) + C_2\sin(\omega_0\,t)\bigr], \qquad \omega_0 = \sqrt{\omega_n^2 - \delta^2}.
\end{equation}
À l'instant du lâcher, l'actionneur est immobile à la position extrême : \(\alpha(0) = \Delta\alpha_{\max}\) et \(\Omega(0) = 0\). De plus, \eqref{eq:meca_libre_omega1} évaluée en \(t=0\) donne \(\dot\Omega(0) = -K_r\,\Delta\alpha_{\max}/J = -\omega_n^2\,\Delta\alpha_{\max}\). Cette relation suppose implicitement \(i(0) = 0\), automatiquement vérifié en circuit ouvert (EX2). Le cas court-circuit (EX3) sera traité en §\ref{sec:laplace} avec \(i_0 = 0\) comme condition initiale au lâcher également (le courant est forcé à zéro tant que la bobine ne bouge pas). Les conditions \(\Omega(0)=0\) et \(\dot\Omega(0) = -\omega_n^2\,\Delta\alpha_{\max}\) imposent \(C_1 = 0\) et \(C_2 = -\omega_n^2\,\Delta\alpha_{\max}/\omega_0\), d'où :
\begin{equation}
\Omega(t) \;=\; -\,\Delta\alpha_{\max}\,\frac{\omega_n^2}{\omega_0}\,e^{-\delta t}\sin(\omega_0\,t).
\label{eq:alpha_t}
\end{equation}

\paragraph{Tension induite de mouvement.}
La bobine étant ouverte, \(i = 0\) et \eqref{eq:elec} se réduit à \(u(t) = K_u\,\Omega(t)\). On a donc :
\begin{equation}
u_{im}(t) \;=\; -\,K_u\,\Delta\alpha_{\max}\,\frac{\omega_n^2}{\omega_0}\,e^{-\delta t}\sin(\omega_0\,t).
\label{eq:uim_t}
\end{equation}
En fonction des paramètres physiques \((J, K_u, C_v, K_r)\), \eqref{eq:uim_t} s'écrit :
\begin{equation}
u_{im}(t) = -\,K_u\,\Delta\alpha_{\max}\,\frac{K_r/J}{\sqrt{K_r/J - C_v^2/(4J^2)}}\,e^{-\frac{C_v}{2J}\,t}\sin\!\left(\sqrt{\tfrac{K_r}{J} - \tfrac{C_v^2}{4J^2}}\;t\right).
\end{equation}

\paragraph{Approximation \(\omega_n \gg \delta\).}
Lorsque l'amortissement est faible, \(\omega_0 \simeq \omega_n\) et \(\omega_n^2/\omega_0 \simeq \omega_n\). L'expression \eqref{eq:uim_t} se simplifie en :
\[
u_{im}(t) \;\approx\; -\,K_u\,\Delta\alpha_{\max}\,\omega_n\,e^{-\delta t}\sin(\omega_n\,t).
\]
La validité de cette simplification sera contrôlée \emph{a posteriori} par le ratio \(\omega_n/\delta\). Pour l'identification numérique de \(K_u\) on garde toutefois la formule exacte \eqref{eq:Ku_ident} de manière à ne pas introduire de biais lié à l'amortissement résiduel (cf.\ §\ref{sec:synthese}).

\paragraph{Identification.}
La pseudo-pulsation \(\omega_0\) se lit directement sur la mesure (\(\omega_0 = 2\pi/T_0\)).

\smallskip
\noindent\emph{Décrément logarithmique.}
L'enveloppe exponentielle \(e^{-\delta t}\) impose, entre deux crêtes successives \(\hat U_1\) (instant \(t_1\)) et \(\hat U_2\) (instant \(t_2 = t_1 + T_0\)) de même signe, le rapport :
\[
\frac{\hat U_1}{\hat U_2} \;=\; \frac{e^{-\delta\,t_1}}{e^{-\delta\,(t_1+T_0)}} \;=\; e^{\delta\,T_0}.
\]
En passant au logarithme, on obtient directement le coefficient d'amortissement :
\begin{equation}
\boxed{\;\delta \;=\; \frac{1}{T_0}\,\ln\!\left(\frac{\hat U_1}{\hat U_2}\right)\;}
\label{eq:delta_decrement}
\end{equation}
Cette mesure se lit directement à l'oscilloscope avec deux curseurs placés sur deux crêtes successives. La pulsation propre non amortie est alors :
\begin{equation}
\omega_n \;=\; \sqrt{\omega_0^2 + \delta^2}.
\label{eq:omegan_ident}
\end{equation}

\smallskip
\noindent\emph{Constante de tension induite.}
Enfin, la première crête de \(u_{im}\), atteinte à \(t \approx T_0/4\), vaut en module :
\[
\hat U_1 \;=\; K_u\,\Delta\alpha_{\max}\,\frac{\omega_n^2}{\omega_0}\,e^{-\delta\,T_0/4},
\]
d'où la formule d'identification de la constante de tension induite :
\begin{equation}
\boxed{\;K_u \;=\; \frac{\hat U_1\;\omega_0}{\omega_n^2\,\Delta\alpha_{\max}\,e^{-\delta\,T_0/4}}\;}
\label{eq:Ku_ident}
\end{equation}
On en déduit également :
\[
\frac{C_v}{J} = 2\,\delta, \qquad \frac{K_r}{J} = \omega_n^2 = \omega_0^2 + \delta^2.
\]

\subsubsection{EX3 -- Bobine court-circuitée : \texorpdfstring{\(J\)}{J}, \texorpdfstring{\(C_v\)}{Cv}, \texorpdfstring{\(K_r\)}{Kr}}

Lorsque les bornes de la bobine sont reliées (\(u = 0\)) et en négligeant \(L\,di/dt\) sous l'hypothèse \(\tau_{el} \ll \tau_{m\acute{e}c}\), \eqref{eq:elec} donne :
\begin{equation}
i(t) \;\approx\; -\,\frac{K_u}{R}\,\Omega(t).
\label{eq:i_short}
\end{equation}
Le couple électromagnétique vaut alors \(T_{em} = K_t\,i = -\,(K_t K_u/R)\,\Omega\). En substituant dans \eqref{eq:meca} et en utilisant \(K_t = K_u\) :
\begin{equation}
J\,\dot\Omega \;=\; -\,K_r\,\alpha \;-\; \!\left(C_v + \frac{K_u^2}{R}\right)\!\Omega.
\label{eq:meca_short_omega1}
\end{equation}
Cette équation a exactement la même forme que \eqref{eq:meca_libre_omega1}, à condition de substituer le frottement effectif :
\[
C_v \;\longrightarrow\; C_v' \;=\; C_v + \frac{K_u^2}{R}.
\]
Toute la résolution menée pour le circuit ouvert se transpose donc directement. En dérivant par rapport au temps, on obtient :
\begin{equation}
\ddot\Omega + 2\delta'\,\dot\Omega + \omega_n^2\,\Omega = 0,
\label{eq:meca_short}
\end{equation}
avec un nouvel amortissement \(\delta'\) qui intègre le freinage électromagnétique :
\begin{equation}
\delta' \;=\; \delta + \frac{K_u^2}{2JR} \;=\; \delta + \frac{K_t\,K_u}{2JR}.
\label{eq:deltap}
\end{equation}
La pulsation propre non amortie \(\omega_n\) reste inchangée puisque \(K_r/J\) ne dépend pas du circuit. L'écart \(\delta' - \delta\) est entièrement attribuable au freinage électromagnétique :
\begin{equation}
\boxed{\;J \;=\; \frac{K_t\,K_u}{2\,R\,(\delta' - \delta)} \;=\; \frac{K_u^2}{2\,R\,(\delta' - \delta)}\;}
\label{eq:J}
\end{equation}
puis :
\[
C_v = 2\,J\,\delta, \qquad K_r = J\,\omega_n^2.
\]

\paragraph{Courant induit.}
Avec les mêmes conditions initiales \(\alpha(0) = \Delta\alpha_{\max}\), \(\Omega(0) = 0\), la solution de \eqref{eq:meca_short} est :
\begin{equation}
\Omega(t) \;=\; -\,\Delta\alpha_{\max}\,\frac{\omega_n^2}{\omega_0'}\,e^{-\delta' t}\sin(\omega_0'\,t), \qquad \omega_0' = \sqrt{\omega_n^2 - \delta'^2}.
\end{equation}
Le courant induit, par \eqref{eq:i_short}, vaut :
\begin{equation}
i(t) \;=\; \frac{K_u\,\Delta\alpha_{\max}}{R}\,\frac{\omega_n^2}{\omega_0'}\,e^{-\delta' t}\sin(\omega_0'\,t).
\label{eq:i_t}
\end{equation}
En fonction des paramètres physiques \((J, K_u, K_t, C_v, K_r, R)\), avec \(\omega_n^2 = K_r/J\) et \(\delta' = (C_v + K_t K_u/R)/(2J)\) :
\begin{equation}
i(t) = \frac{K_u\,\Delta\alpha_{\max}}{R}\,\frac{K_r/J}{\sqrt{K_r/J - \delta'^2}}\,e^{-\delta' t}\sin\!\left(\sqrt{\tfrac{K_r}{J} - \delta'^2}\;t\right).
\end{equation}

\paragraph{Validité de \(\omega_n \gg \delta'\).}
Le freinage électromagnétique étant nettement supérieur au frottement visqueux mécanique, l'approximation \(\omega_n \gg \delta'\) n'est pas garantie en court-circuit. Le ratio \(\omega_n/\delta'\) sera donc évalué expérimentalement.

\paragraph{Identification.}
La même méthode du décrément logarithmique s'applique sur le courant. Avec deux crêtes successives de même signe \(\hat I_1\) et \(\hat I_2\) séparées de \(T_0'\) :
\begin{equation}
\delta' \;=\; \frac{1}{T_0'}\,\ln\!\left(\frac{\hat I_1}{\hat I_2}\right), \qquad \omega_0' = \frac{2\pi}{T_0'}.
\label{eq:deltap_decrement}
\end{equation}
On vérifie que \(\omega_n^2 = \omega_0'^2 + \delta'^2\) coïncide avec la valeur extraite en EX2 (la raideur \(K_r\) ne dépendant pas du circuit).

\subsection{Détermination du nombre de spires}
\label{sec:nbre_spires}

Une fois \(K_u\) identifié et connaissant l'induction dans l'entrefer \(B\) ainsi que les rayons \(R_e\) et \(R_i\) de la bobine, la relation géométrique \eqref{eq:Ku_geom} permet de remonter au nombre de spires :
\begin{equation}
\boxed{\;N \;=\; \frac{K_u}{B\,(R_e^2 - R_i^2)}\;}
\label{eq:N_spires}
\end{equation}

\subsection{Transformées de Laplace sans simplification}
\label{sec:laplace}

Pour répondre au point §2.2.4-4 de la donnée, on écrit d'abord les équations générales \eqref{eq:elec} et \eqref{eq:meca} dans le domaine de Laplace, sans imposer \(\tau_{el} \ll \tau_{m\acute{e}c}\), et en tenant compte des conditions initiales \(\alpha(0) = \alpha_0\), \(\dot\alpha(0) = \Omega_0\), \(i(0) = i_0\) :
\[
\begin{aligned}
U(s) &= (R + L\,s)\,I(s) - L\,i_0 + K_u\,\bigl[s\,\alpha(s) - \alpha_0\bigr], \\
\bigl(J s^2 + C_v s + K_r\bigr)\alpha(s) &= K_t\,I(s) + J\,\bigl(s\,\alpha_0 + \Omega_0\bigr) + C_v\,\alpha_0.
\end{aligned}
\]

\paragraph{Bobine ouverte.}
Dans l'essai du §2.2.2.1, \(i(t)=0\). En posant \(\Omega=\dot\alpha\), l'équation différentielle obtenue plus haut est :
\[
\ddot\Omega + 2\delta\,\dot\Omega + \omega_n^2\,\Omega = 0,
\qquad
\delta=\frac{C_v}{2J},\quad \omega_n^2=\frac{K_r}{J}.
\]
Sa transformée de Laplace, avec \(\Omega(0)=\Omega_0\) et \(\dot\Omega(0)=\dot\Omega_0\), vaut :
\begin{equation}
\bigl(s^2+2\delta s+\omega_n^2\bigr)\Omega(s)
= s\Omega_0+\dot\Omega_0+2\delta\Omega_0.
\label{eq:laplace_open_general}
\end{equation}
Pour le lâcher depuis la butée, \(\alpha(0)=\alpha_0=\Delta\alpha_{\max}\), \(\Omega_0=0\) et \(\dot\Omega_0=-\omega_n^2\alpha_0\), d'où :
\begin{equation}
\Omega(s)=\frac{-\omega_n^2\alpha_0}{s^2+2\delta s+\omega_n^2},
\qquad
U_{im}(s)=K_u\,\Omega(s).
\label{eq:laplace_open_release}
\end{equation}

\paragraph{Bobine court-circuitée sans négliger l'inductance.}
Dans l'essai du §2.2.3.1, \(u(t)=0\), mais on garde le terme \(L\,di/dt\). Les équations couplées sont donc :
\[
L\frac{di}{dt}+R\,i+K_u\,\Omega=0,
\qquad
J\dot\Omega+C_v\Omega+K_r\alpha=K_t i.
\]
En éliminant \(i\) et \(\alpha\), on obtient une équation différentielle d'ordre 3 pour la vitesse :
\begin{equation}
LJ\,\Omega^{(3)}
+(LC_v+RJ)\,\ddot\Omega
+(LK_r+RC_v+K_tK_u)\,\dot\Omega
+RK_r\,\Omega = 0.
\label{eq:short_full_ode}
\end{equation}
Avec \(\Omega(0)=\Omega_0\), \(\dot\Omega(0)=\dot\Omega_0\) et \(\ddot\Omega(0)=\ddot\Omega_0\), sa transformée de Laplace s'écrit :
\begin{equation}
\begin{aligned}
&\Bigl[LJ\,s^3+(LC_v+RJ)s^2+(LK_r+RC_v+K_tK_u)s+RK_r\Bigr]\Omega(s) \\
&= LJ\bigl(s^2\Omega_0+s\dot\Omega_0+\ddot\Omega_0\bigr)
+(LC_v+RJ)\bigl(s\Omega_0+\dot\Omega_0\bigr)
+(LK_r+RC_v+K_tK_u)\Omega_0.
\end{aligned}
\label{eq:laplace_short_full}
\end{equation}
Pour un lâcher depuis la butée au repos, avec \(i_0=0\), on a \(\Omega_0=0\), \(\dot\Omega_0=-K_r\alpha_0/J\) et, par l'équation électrique, \(\dot i(0)=0\), donc \(\ddot\Omega_0=C_vK_r\alpha_0/J^2\).

En conditions initiales nulles (\(i_0 = \alpha_0 = \Omega_0 = 0\)), la fonction de transfert globale \(\alpha(s)/U(s)\) s'obtient par élimination de \(I(s)\) :
\begin{equation}
\frac{\alpha(s)}{U(s)} \;=\; \frac{K_t}{(R + L s)\,(J s^2 + C_v s + K_r) + K_u K_t\,s}.
\label{eq:FT}
\end{equation}
Le dénominateur est un polynôme du \emph{troisième} ordre en \(s\) ; il se factorise en un mode électrique rapide à \(s \simeq -R/L\) et une dynamique mécanique d'ordre 2 à \(s \simeq -\delta' \pm j\omega_0'\) lorsque \(\tau_{el} \ll \tau_{m\acute{e}c}\), retrouvant ainsi le modèle d'ordre 2 utilisé pour l'identification.

\subsection{Schéma fonctionnel (bonus §2.2.5)}

\begin{figure}[H]
\centering
\resizebox{\textwidth}{!}{%
\begin{tikzpicture}[>=Stealth, node distance=1.5cm and 1.6cm,
   block/.style={draw, thick, rounded corners=2pt, minimum height=1cm, minimum width=2.4cm, align=center},
   sum/.style={draw, circle, thick, inner sep=1pt, minimum size=5mm}]

\node[coordinate] (in) {};
\node[sum, right=1cm of in] (sum1) {};
\node[block, right=of sum1] (elec) {$\dfrac{1}{R + L\,s}$};
\node[block, right=of elec] (Kt) {$K_t$};
\node[sum, right=of Kt] (sum2) {};
\node[block, right=of sum2] (mec) {$\dfrac{1}{J s^2 + C_v s + K_r}$};
\node[coordinate, right=1cm of mec] (out) {};
\node[block, below=1.4cm of Kt] (Ku) {$K_u\,s$};

\draw[->] (in) -- (sum1) node[pos=0,above]{$U(s)$};
\draw[->] (sum1) -- (elec) node[pos=0.5,above]{};
\draw[->] (elec) -- (Kt) node[pos=0.5,above]{$I(s)$};
\draw[->] (Kt) -- (sum2);
\draw[->] (sum2) -- (mec);
\draw[->] (mec) -- (out) node[pos=1,above]{$\alpha(s)$};

% feedback fcem
\draw (mec) |- (Ku);
\draw[->] (Ku) -| (sum1);
\node[below right=0.05 and 0.1 of sum1] {\scriptsize $-$};

% perturbation de couple (option = 0 ici)
\node[above=0.4 of sum2] {\scriptsize $T_{ext}=0$};

\end{tikzpicture}%
}
\caption{Schéma fonctionnel complet U $\to \alpha$ (bonus §2.2.5)}
\label{fig:schema_fonctionnel}
\end{figure}
